\documentclass[12pt]{article}
\usepackage[brazil]{babel}
\usepackage[utf8]{inputenc}
\usepackage[T1]{fontenc}
\usepackage{amsmath, amssymb}
\usepackage{geometry}
\usepackage{graphicx}
\usepackage{xcolor}

\definecolor{maincolor}{RGB}{0,100,50}
\usepackage{fancyhdr}
\definecolor{forestgreen}{RGB}{0,100,50}
\usepackage[most]{tcolorbox}
\newtcolorbox{block}[1]{%
  colback=green!5,
  colframe=maincolor,
  coltitle=white,
  title=#1,
  fonttitle=\bfseries,
  boxrule=1.5pt,
  arc=3mm,
  left=6pt,
  right=6pt,
  top=6pt,
  bottom=6pt
}

\geometry{margin=2.5cm}
\pagestyle{fancy}
\fancyhf{} % limpa cabeçalho e rodapé

% Linha do cabeçalho e rodapé
\renewcommand{\headrulewidth}{2pt}
\renewcommand{\footrulewidth}{2pt}

\renewcommand{\headrule}{\hbox to\headwidth{%
  \color{forestgreen}\leaders\hrule height \headrulewidth\hfill}}

\renewcommand{\footrule}{\hbox to\headwidth{%
  \color{forestgreen}\leaders\hrule height \footrulewidth\hfill}}
% Ajuste de altura do cabeçalho
\setlength{\headheight}{32pt}

% Cabeçalho com imagem 
\fancyhead[L]{%
  \includegraphics[height=1.2cm]{Universidade_de_Rio_Verde.jpg}
}

\fancyhead[R]{%
  \raggedleft
  \textcolor{forestgreen}{%
    \small
    \textbf{Lista de Exercícios — Provas Matemáticas}\\
    Me.~Sergio Souza Novak\\
    \textit{Lógica para Engenharia de Software}
  }
}

% Rodapé
\fancyfoot[C]{%
  \textcolor{forestgreen}{\thepage}
}

\begin{document}
% \maketitle
\vspace*{0.2cm}
 
\section*{Instruções}
\begin{block}{Utilize as definições formais de números pares e ímpares}
Um inteiro é par se pode ser escrito na forma \( n = 2k \), com \( k \in \mathbb{Z} \).


\[
    n \in \mathbb{Z} \text{ é par} \iff \exists k \in \mathbb{Z} \text{ tal que } n = 2k,
    \]
    \[
        n \in \mathbb{Z} \text{ é ímpar} \iff \exists k \in \mathbb{Z} \text{ tal que } n = 2k + 1.
        \]
    \end{block}
Justifique rigorosamente todas as afirmações.

\section*{Exercícios}

\begin{enumerate}

\item Demonstre que a soma de dois números pares é sempre um número par.
\begin{block}{Resposta}
Sejam $a$ e $b$ dois números inteiros pares. Pela definição de número par, existem
$k,m \in \mathbb{Z}$ tais que $a = 2k$ e $b = 2m$. Assim, $ a+b = 2k + 2m = 2(k+m)$
onde $k+m \in \mathbb{Z}$. Logo, $a+b$ é par.
\end{block}


\item Demonstre que a soma de dois números ímpares é sempre um número par.
\begin{block}{Resposta}
Sejam $a$ e $b$ dois números inteiros ímpares. Existem $k,m \in \mathbb{Z}$ tais que
$a = 2k+1$ e $b = 2m+1$. Portanto,
$a+b = 2k+1 + 2m+1 = 2(k+m+1)$,
com $k+m+1 \in \mathbb{Z}$. Assim, a soma é par.
\end{block}
\item Demonstre que a soma de um número par com um número ímpar é sempre um número ímpar.

\item Prove que o produto de dois números inteiros ímpares é sempre ímpar.

\item Demonstre que o quadrado de qualquer número inteiro ímpar é ímpar.

\item Demonstre que, para qualquer inteiro \( n \), exatamente um dos números \( n \) ou \( n+1 \) é par.

\item Prove que, para todo inteiro \( n \), o número
\[
3n + 1
\]
possui paridade oposta à de \( n \).

\item Demonstre que a soma de três números inteiros consecutivos é sempre ímpar.
\newpage
\end{enumerate}


\begin{block}{3. Soma de um inteiro par com um inteiro ímpar}
Seja $a$ um inteiro par e $b$ um inteiro ímpar. Existem $k,m \in \mathbb{Z}$ tais que
$a = 2k$ e $b = 2m+1$. Logo,
\[
a+b = 2k + 2m + 1 = 2(k+m)+1,
\]
onde $k+m \in \mathbb{Z}$. Portanto, a soma é ímpar.
\end{block}

\begin{block}{4. Produto de dois inteiros ímpares}
Sejam $a$ e $b$ dois inteiros ímpares. Existem $k,m \in \mathbb{Z}$ tais que
$a = 2k+1$ e $b = 2m+1$. O produto é
\[
ab = (2k+1)(2m+1) = 2(2km+k+m)+1,
\]
com $2km+k+m \in \mathbb{Z}$. Logo, $ab$ é ímpar.
\end{block}

\begin{block}{5. Quadrado de um inteiro ímpar}
Seja $n$ um inteiro ímpar. Existe $k \in \mathbb{Z}$ tal que $n = 2k+1$. Assim,
\[
n^2 = (2k+1)^2 = 2(2k^2+2k)+1,
\]
onde $2k^2+2k \in \mathbb{Z}$. Portanto, o quadrado de um inteiro ímpar é ímpar.
\end{block}

\begin{block}{6. Paridade de $n$ e $n+1$}
Seja $n \in \mathbb{Z}$. Se $n$ é par, então $n=2k$ e $n+1=2k+1$, que é ímpar.
Se $n$ é ímpar, então $n=2k+1$ e $n+1=2(k+1)$, que é par.
Assim, exatamente um dos números $n$ ou $n+1$ é par.
\end{block}

\begin{block}{7. Paridade de $3n+1$}
Seja $n \in \mathbb{Z}$. Se $n=2k$, então
\[
3n+1 = 6k+1 = 2(3k)+1,
\]
que é ímpar. Se $n=2k+1$, então
\[
3n+1 = 6k+4 = 2(3k+2),
\]
que é par. Logo, $3n+1$ tem paridade oposta à de $n$.
\end{block}

\begin{block}{8. Soma de três inteiros consecutivos}
Sejam três inteiros consecutivos $n-1$, $n$ e $n+1$. A soma é
\[
(n-1)+n+(n+1)=3n.
\]
Se $n$ é par, então $3n$ é par; se $n$ é ímpar, então $3n$ é ímpar.
Em ambos os casos, a soma dos três inteiros consecutivos é ímpar.
\end{block}

\end{document}
